\documentclass[12pt,a4paper]{report}

\usepackage[margin=1in]{geometry}
\usepackage{xcolor}
\usepackage{sectsty}
\usepackage{fancyhdr}

\title{Introduction to LaTeX}
\author{Saarvik Singh Suryavanshi}
\date{}

\chapterfont{\color{red}\bfseries\Huge}
\sectionfont{\color{magenta}\itshape\Large}
\subsectionfont{\color{orange}\bfseries}

\pagestyle{fancy}
\fancyhf{}

\fancyhead[L]{Saarvik Singh Suryavanshi}
\fancyhead[R]{Technical Writing}

\fancyfoot[C]{\thepage}

\renewcommand{\headrulewidth}{0.4pt}
\renewcommand{\footrulewidth}{0.4pt}

\fancypagestyle{plain}{
	\fancyhf{}
	\fancyhead[L]{Saarvik Singh Suryavanshi}
	\fancyhead[R]{Technical Writing}
	\fancyfoot[C]{\thepage}
	\renewcommand{\headrulewidth}{0.4pt}
	\renewcommand{\footrulewidth}{0.4pt}
}

\begin{document}
	
	\begin{titlepage}
		\thispagestyle{empty}
		
		\centering
		
		\vspace*{3cm}
		
		{\Huge\bfseries Introduction to LaTeX\par}
		
		\vspace{2cm}
		
		{\Large Saarvik Singh Suryavanshi\par}
		
		\vspace{0.5cm}
		
		2025UCP1681
		
		\vspace{1.5cm}
		
		\textbf{Course:} Technical Writing
		
		\vspace{0.5cm}
		
		\textbf{Course Coordinator:} Neeta Nain
		
		\vfill
		
		{\large \par}
		
	\end{titlepage}
	
	\pagestyle{fancy}
	
	\begin{abstract}
		
		This report introduces LaTeX and explains how it can be
		used to prepare well-formatted documents. It covers the
		basic document structure, commonly used commands,
		mathematical expressions, and the use of tables and figures.
		The aim is to understand the basic features of LaTeX and
		how they help in preparing academic documents.
		
	\end{abstract}
	
	\section{Introduction}
	
	LaTeX is a document preparation system commonly used for
	writing academic papers, reports, and technical documents.
	It is especially useful when a document contains
	mathematical expressions, tables, references, or figures.
	
	Unlike a normal word processor, LaTeX uses commands to
	describe the structure and formatting of a document.
	This makes it easier to maintain consistent formatting
	throughout a document.
	
	\chapter{Introduction to LaTeX}
	
	\section{What is LaTeX?}
	
	LaTeX is built on the TeX typesetting system. It allows
	users to focus on the content while the system handles
	most of the formatting.
	
	Some common uses of LaTeX are:
	
	\begin{itemize}
		\item Writing technical and research papers.
		\item Preparing project reports and assignments.
		\item Creating documents with mathematical equations.
		\item Preparing presentations.
	\end{itemize}
	
	\section{Advantages of LaTeX}
	
	LaTeX provides several useful features:
	
	\begin{itemize}
		\item Consistent formatting throughout the document.
		\item Good support for mathematical equations.
		\item Automatic numbering of sections, tables, and figures.
		\item Easy management of references and citations.
	\end{itemize}
	
	\chapter{Basic Document Structure}
	
	\section{Document Class}
	
	Every LaTeX document begins with the
	\texttt{\textbackslash documentclass} command.
	
	It specifies the type of document being prepared.
	Some commonly used document classes are:
	
	\begin{itemize}
		\item article
		\item report
		\item book
		\item beamer
	\end{itemize}
	
	\section{Basic Commands}
	
	A LaTeX document generally contains a preamble and
	a document environment.
	
	The preamble contains the document class and packages.
	The actual content is written between
	\texttt{\textbackslash begin\{document\}} and
	\texttt{\textbackslash end\{document\}}.
	
	\subsection{Sections and Subsections}
	
	Sections are used to divide a document into smaller
	parts. Subsections help organize the content further.
	
	\chapter{Mathematics, Tables and Figures}
	
	\section{Mathematical Expressions}
	
	LaTeX is useful for writing mathematical expressions.
	For example, the quadratic formula can be written as:
	
	\[
	x = \frac{-b \pm \sqrt{b^2-4ac}}{2a}
	\]
	
	Equations can also be numbered and referenced in
	a document.
	
	\section{Tables and Figures}
	
	Tables are used to present information in rows and
	columns. Figures can be used to include images,
	graphs, and diagrams.
	
	Both can be given captions and labels, which make
	them easier to identify and reference in a report.
	
	\chapter{Conclusion}
	
	LaTeX is a useful tool for preparing academic and
	technical documents. Its structured approach makes
	formatting easier to manage, especially in longer
	documents.
	
	Learning its basic commands, environments, and
	packages provides a good starting point for creating
	well-organized reports and presentations.
	
\end{document}